\documentclass[10pt,a4paper]{amsart}
\usepackage{amsmath,amssymb,amsthm,graphicx}
\usepackage{verbatim}
\newenvironment{reference}{\comment}{\endcomment}
%\newenvironment{reference}{}{}
\newenvironment{slogan}{\comment}{\endcomment}
\newenvironment{history}{\comment}{\endcomment}
% For commutative diagrams we use Xy-pic
\usepackage[all]{xy}

% We use 2cell for 2-commutative diagrams.
\xyoption{2cell}
% Theorem environments.
%
\theoremstyle{plain}
\newtheorem{theorem}[subsection]{Theorem}
\newtheorem{proposition}[subsection]{Proposition}
\newtheorem{lemma}[subsection]{Lemma}

\theoremstyle{definition}
\newtheorem{definition}[subsection]{Definition}
\newtheorem{example}[subsection]{Example}
\newtheorem{exercise}[subsection]{Exercise}
\newtheorem{situation}[subsection]{Situation}

\theoremstyle{remark}
\newtheorem{remark}[subsection]{Remark}
\newtheorem{remarks}[subsection]{Remarks}

\numberwithin{equation}{subsection}

% Macros
%
\def\lim{\mathop{\mathrm{lim}}\nolimits}
\def\colim{\mathop{\mathrm{colim}}\nolimits}
\def\Spec{\mathop{\mathrm{Spec}}}
\def\Hom{\mathop{\mathrm{Hom}}\nolimits}
\def\Ext{\mathop{\mathrm{Ext}}\nolimits}
\def\SheafHom{\mathop{\mathcal{H}\!\mathit{om}}\nolimits}
\def\SheafExt{\mathop{\mathcal{E}\!\mathit{xt}}\nolimits}
\def\Sch{\mathit{Sch}}
\def\Mor{\mathop{\mathrm{Mor}}\nolimits}
\def\Ob{\mathop{\mathrm{Ob}}\nolimits}
\def\Sh{\mathop{\mathit{Sh}}\nolimits}
\def\NL{\mathop{N\!L}\nolimits}
\def\CH{\mathop{\mathrm{CH}}\nolimits}
\def\proetale{{pro\text{-}\acute{e}tale}}
\def\etale{{\acute{e}tale}}
\def\QCoh{\mathit{QCoh}}
\def\Ker{\mathop{\mathrm{Ker}}}
\def\Im{\mathop{\mathrm{Im}}}
\def\Coker{\mathop{\mathrm{Coker}}}
\def\Coim{\mathop{\mathrm{Coim}}}

% Boxtimes
%
\DeclareMathSymbol{\boxtimes}{\mathbin}{AMSa}{"02}

%
% Macros for moduli stacks/spaces
%
\def\QCohstack{\mathcal{QC}\!\mathit{oh}}
\def\Cohstack{\mathcal{C}\!\mathit{oh}}
\def\Spacesstack{\mathcal{S}\!\mathit{paces}}
\def\Quotfunctor{\mathrm{Quot}}
\def\Hilbfunctor{\mathrm{Hilb}}
\def\Curvesstack{\mathcal{C}\!\mathit{urves}}
\def\Polarizedstack{\mathcal{P}\!\mathit{olarized}}
\def\Complexesstack{\mathcal{C}\!\mathit{omplexes}}
% \Pic is the operator that assigns to X its picard group, usage \Pic(X)
% \Picardstack_{X/B} denotes the Picard stack of X over B
% \Picardfunctor_{X/B} denotes the Picard functor of X over B
\def\Pic{\mathop{\mathrm{Pic}}\nolimits}
\def\Picardstack{\mathcal{P}\!\mathit{ic}}
\def\Picardfunctor{\mathrm{Pic}}
\def\Deformationcategory{\mathcal{D}\!\mathit{ef}}
\usepackage{cleveref}
\crefname{originalremark}{original remark}{original remarks}
\Crefname{originalremark}{Original remark}{Original remarks}
\makeatletter
\newlabel{cohomology-remark-support-cup-product}{{0G75}{source}{}{sourcepointer-1}{}}
\newlabel{cohomology-remark-support-cup-product@cref}{{[originalremark][0][]0G75}{[1][0][]source}{}{}{}}
\makeatother
\crefname{originalsection}{original section}{original sections}
\Crefname{originalsection}{Original section}{Original sections}
\makeatletter
\newlabel{cohomology-section-cup-product}{{0FKU}{source}{}{sourcepointer-2}{}}
\newlabel{cohomology-section-cup-product@cref}{{[originalsection][0][]0FKU}{[1][0][]source}{}{}{}}
\makeatother
\crefname{originallemma}{original lemma}{original lemmas}
\Crefname{originallemma}{Original lemma}{Original lemmas}
\makeatletter
\newlabel{lemma-gysin-global}{{0G86}{source}{}{sourcepointer-3}{}}
\newlabel{lemma-gysin-global@cref}{{[originallemma][0][]0G86}{[1][0][]source}{}{}{}}
\makeatother
\crefname{originallemma}{original lemma}{original lemmas}
\Crefname{originallemma}{Original lemma}{Original lemmas}
\makeatletter
\newlabel{perfect-lemma-supported-trivial-vanishing}{{0G7N}{source}{}{sourcepointer-4}{}}
\newlabel{perfect-lemma-supported-trivial-vanishing@cref}{{[originallemma][0][]0G7N}{[1][0][]source}{}{}{}}
\makeatother
\title{The Gysin map on a restricted class is the cup product of the pulled back class with the Gysin class of the unit}
\author{The Stacks Project Authors}
\date{stacks-project:derham.tex:local-snapshot-20261012; GNU FDL 1.2 or later}
\begin{document}
\maketitle
\noindent\textit{Editorial note.} All delivered mathematical text is the author's own, copied from the pinned Stacks Project chapter file derham.tex; no statement or proof step was rewritten, reordered or removed.\par\medskip
\noindent\textit{Editorial note.} Preamble provenance: the pinned chapter file opens with \textbackslash{}input\{preamble\} and carries no preamble text of its own, so the project's single Stacks-wide preamble staged at sources/stacks-derham/preamble.tex is excerpted verbatim as three spans, lines 8-12 (verbatim package and the project's own reference/slogan/history display environments), lines 14-18 (Xy-pic setup) and lines 154-217 (theorem declarations and mathematical macros). The wrapper class and the amsmath/amssymb/amsthm/graphicx base packages are supplied by the assembler. No macro was invented or redefined.\par\medskip
\noindent\textit{Editorial note.} Same-chapter pointers are located at their \textbackslash{}label lines inside the pinned chapter file; cross-chapter pointers are located in the project's own tag list staged in the snapshot (sources/stacks-chow/tags.txt), at the verbatim lines carrying their canonical tags.\par\medskip
\noindent\textit{Allowed standard prerequisites (editorial).}\begin{itemize}
\item The global Gysin setup and the maps on local cohomology. Same chapter derham.tex, label lemma-gysin-global (cited at statement line 5521). It fixes the situation of \$i : Z \textbackslash{}to X\$ and provides the maps \$\textbackslash{}gamma\textasciicircum{}p : \textbackslash{}Omega\textasciicircum{}p\_\{Z/S\} \textbackslash{}to \textbackslash{}mathcal\{H\}\textasciicircum{}c\_Z(\textbackslash{}Omega\textasciicircum{}\{p+c\}\_\{X/S\})\$ whose interaction with the cup product is the assertion.
\item The cup product for local cohomology with support, and the general cup product formalism. Cohomology, remark cohomology-remark-support-cup-product (invoked at proof line 5551) and Section \textbackslash{}ref\{cohomology-section-cup-product\} (invoked at proof line 5577). The first supplies the canonical maps from the tensor product with \$R\textbackslash{}mathcal\{H\}\_Z(\textbackslash{}Omega\textasciicircum{}c\_\{X/S\})\$ into the local cohomology of the tensor product, the second lets the claim be reduced to the commutativity of a single diagram.
\item Vanishing of the local cohomology sheaves above the codimension. Derived Categories of Schemes, label perfect-lemma-supported-trivial-vanishing (invoked at proof line 5560). It shows that \$R\textbackslash{}mathcal\{H\}\_Z(\textbackslash{}Omega\textasciicircum{}j\_\{X/S\})\$ has no cohomology above degree \$c\$, so passing to degree \$c\$ cohomology sheaves is legitimate.
\end{itemize}\medskip
\ifdefined\hypertarget\hypertarget{sourcepointer-1}{}\fi
\noindent\textit{Source reference.} Original remark 0G75: sources/stacks-chow/tags.txt, source lines 20220--20220. Navigation pointer for \textbackslash{}ref\{cohomology-remark-support-cup-product\} used inside the counted span (line 5551); the argument that consumes it is retained in the span.\par\medskip
\ifdefined\hypertarget\hypertarget{sourcepointer-2}{}\fi
\noindent\textit{Source reference.} Original section 0FKU: sources/stacks-chow/tags.txt, source lines 19474--19474. Navigation pointer for \textbackslash{}ref\{cohomology-section-cup-product\} used inside the counted span (line 5577); the argument that consumes it is retained in the span.\par\medskip
\ifdefined\hypertarget\hypertarget{sourcepointer-3}{}\fi
\noindent\textit{Source reference.} Original lemma 0G86: sources/stacks-derham/pinned-original-source.tex, source lines 5473--5473. Navigation pointer for \textbackslash{}ref\{lemma-gysin-global\} used inside the counted span (line 5521); the argument that consumes it is retained in the span.\par\medskip
\ifdefined\hypertarget\hypertarget{sourcepointer-4}{}\fi
\noindent\textit{Source reference.} Original lemma 0G7N: sources/stacks-chow/tags.txt, source lines 20238--20238. Navigation pointer for \textbackslash{}ref\{perfect-lemma-supported-trivial-vanishing\} used inside the counted span (line 5561); the argument that consumes it is retained in the span.\par\medskip
\begin{lemma}
\label{lemma-gysin-projection}
Let $X \to S$ and $i : Z \to X$ be as in Lemma \ref{lemma-gysin-global}.
Given $\alpha \in H^q(X, \Omega^p_{X/S})$ we have
$\gamma^p(\alpha|_Z) = i^{-1}\alpha \wedge \gamma^0(1)$ in
$H^q(Z, \mathcal{H}^c_Z(\Omega^{p + c}_{X/S}))$.
Please see proof for notation.
\end{lemma}

\begin{proof}
The restriction $\alpha|_Z$ is the element of $H^q(Z, \Omega^p_{Z/S})$
given by functoriality for Hodge cohomology. Applying functoriality
for cohomology using
$\gamma^p : \Omega^p_{Z/S} \to \mathcal{H}^c_Z(\Omega^{p + c}_{X/S})$
we get get $\gamma^p(\alpha|_Z)$ in
$H^q(Z, \mathcal{H}^c_Z(\Omega^{p + c}_{X/S}))$.
This explains the left hand side of the formula.

\medskip\noindent
To explain the right hand side, we first pullback by the map
of ringed spaces $i : (Z, i^{-1}\mathcal{O}_X) \to (X, \mathcal{O}_X)$
to get the element $i^{-1}\alpha \in H^q(Z, i^{-1}\Omega^p_{X/S})$.
Let $\gamma^0(1) \in H^0(Z, \mathcal{H}_Z^c(\Omega^c_{X/S}))$
be the image of $1 \in H^0(Z, \mathcal{O}_Z) = H^0(Z, \Omega^0_{Z/S})$
by $\gamma^0$. Using cup product we obtain an element
$$
i^{-1}\alpha \cup \gamma^0(1)
\in
H^{q + c}(Z, 
i^{-1}\Omega^p_{X/S} \otimes_{i^{-1}\mathcal{O}_X}
\mathcal{H}^c_Z(\Omega^c_{X/S}))
$$
Using Cohomology, Remark \ref{cohomology-remark-support-cup-product}
and wedge product there are canonical maps
$$
i^{-1}\Omega^p_{X/S} \otimes_{i^{-1}\mathcal{O}_X}^\mathbf{L}
R\mathcal{H}_Z(\Omega^c_{X/S}) \to
R\mathcal{H}_Z(\Omega^p_{X/S} \otimes_{\mathcal{O}_X}^\mathbf{L}
\Omega^c_{X/S}) \to
R\mathcal{H}_Z(\Omega^{p + c}_{X/S})
$$
By Derived Categories of Schemes, Lemma
\ref{perfect-lemma-supported-trivial-vanishing}
the objects $R\mathcal{H}_Z(\Omega^j_{X/S})$ have vanishing
cohomology sheaves in degrees $> c$. Hence on cohomology
sheaves in degree $c$ we obtain a map
$$
i^{-1}\Omega^p_{X/S} \otimes_{i^{-1}\mathcal{O}_X}
\mathcal{H}^c_Z(\Omega^c_{X/S}) \longrightarrow
\mathcal{H}^c_Z(\Omega^{p + c}_{X/S})
$$
The expression $i^{-1}\alpha \wedge \gamma^0(1)$ is the image
of the cup product $i^{-1}\alpha \cup \gamma^0(1)$ by the
functoriality of cohomology.

\medskip\noindent
Having explained the content of the formula in this manner, by
general properties of cup products
(Cohomology, Section \ref{cohomology-section-cup-product}),
it now suffices to prove that the diagram
$$
\xymatrix{
i^{-1}\Omega^p_X \otimes \Omega^0_Z \ar[rr]_{\text{id} \otimes \gamma^0}
\ar[d] & &
i^{-1}\Omega^p_X \otimes \mathcal{H}^c_Z(\Omega^c_X) \ar[d]^\wedge \\
\Omega^p_Z \otimes \Omega^0_Z \ar[r]^\wedge &
\Omega^p_Z \ar[r]^{\gamma^p} &
\mathcal{H}^c_Z(\Omega^{p + c}_X)
}
$$
is commutative in the category of sheaves on $Z$ (with obvious abuse of
notation). This boils down to a simple computation for the maps
$\gamma^j_{f_1, \ldots, f_c}$ which we omit; in fact these maps
are chosen exactly such that this works and such that $1$ maps to
$\frac{\text{d}f_1 \wedge \ldots \wedge \text{d}f_c}{f_1 \ldots f_c}$.
\end{proof}

\end{document}
